\documentclass[10pt,a4paper]{amsart}
\usepackage[margin=19mm]{geometry}
\usepackage{amsmath,amssymb,mathtools}
\usepackage[scr=rsfs,cal=euler]{mathalfa}
\usepackage{xfrac}
\usepackage[unicode,hidelinks,bookmarks=false]{hyperref}
\newcommand{\ie}{i.e.\ }
\newcommand{\cf}{cf.\ }
\newcommand{\resp}{respectively\ }
\newcommand{\Subsection}{\subsection}
\providecommand{\nobreakdash}{\nobreak\hbox{-}}
\setlength{\emergencystretch}{2em}
\multlinegap0pt
\usepackage{mathtools}
\usepackage{amssymb}
\usepackage[shortlabels]{enumitem}
\usepackage{xcolor}
\newtheorem{proposition}{Proposition}
\newcommand{\Cart}{\operatorname{Cart}}
\newcommand{\Frob}{\operatorname{Frob}}
\DeclareMathOperator{\Fun}{\mathrm{Fun}}
\newcommand{\Map}{\operatorname{Map}}
\newcommand{\cC}{\mathcal{C}}
\newcommand{\ev}{\mathrm{ev}}
\newcommand{\op}{{\operatorname{op}}}
\title{Frobenius--Cartier duality via lax equalizers}
\author{Fei Ren}
\date{Source: arXiv:2609.10140v1 (CC BY 4.0)}
\begin{document}
\maketitle

\noindent\emph{Source and licence.} Frobenius--Cartier duality via lax equalizers. Version arXiv:2609.10140v1 (CC BY 4.0), distributed by arXiv under the Creative Commons Attribution 4.0 International licence, http://creativecommons.org/licenses/by/4.0/, as recorded in the arXiv licence metadata for this exact version and shown on the abstract page https://arxiv.org/abs/2609.10140v1. Full licence notice: \texttt{licenses/arxiv-2609.10140-CC-BY-4.0.txt}.\par\smallskip

\noindent\textit{Editorial scope.} Counted unit: the article's proposition prop:abstract-hom (source0.tex lines 915--943). The article's complete original proof of it is delivered with it (source0.tex lines 944--990, one \texttt{proof} environment of 47 lines). No mathematics is written, repaired or re-derived: the statement and the proof are the article's own lines, in the article's own notation, and no retained line is substituted or re-flowed.  No further source-local context is needed: every cross-reference of every retained line resolves inside the two delivered spans.  Nothing was omitted from any retained span, so the package carries no omission record.  The retained lines cite 2 works, whose entries are transcribed verbatim from the article's own bibliography source in the e-print and delivered with short editorial labels recorded in the item's reference mapping.  No retained line contains a figure environment, an image-inclusion command or drawing code, so no image is delivered and the package declares no asset dependency.  Editorial formatting only, in addition: an AMS \texttt{amsart} preamble that restates the article's own declarations and macros used by the retained lines, namely the environments \texttt{proposition} on separate counters, together with the standard packages those lines need.  The retained environments print in this document as Proposition 1 in order of appearance, whereas the article prints the same items with the numbering of its own full document.  The complete native source of the article as distributed in the arXiv e-print for this exact version is bundled unchanged as \texttt{sources/209878/source0.tex}.
\begin{proposition}[Internal Hom pairing]\label{prop:abstract-hom}
Let $\cC$ be a closed monoidal $\infty$-category, and let
$C\mapsto B^C$ denote the right adjoint to $A\mapsto A\otimes B$,
as in \cite[Definition 4.1.1.15; Remark 4.1.1.16]{HA}.
Let $G:\cC\to\cC$ carry a lax monoidal $\infty$-functor structure
in the sense of \cite[Definition 2.1.2.7; Example 2.2.6.10]{HA},
with structure map $\mu_{A,B}:GA\otimes GB\to G(A\otimes B)$.
Then:
\begin{enumerate}
\item
There is a natural transformation
\[
 \chi_{B,C}:G B^C\longrightarrow (GB)^{GC}
\]
in $\Fun(\cC^{\op}\times\cC,\cC)$.
\item
It induces a functor
\[
 \Frob(\cC,G)^{\op}\times\Cart(\cC,G)
       \longrightarrow\Cart(\cC,G)
\]
whose value on $(B,\tau:B\to GB)$ and $(C,\kappa:GC\to C)$
has underlying object $B^C$ and structure map
\begin{equation}\label{eq:abstract-hom-structure}
 G B^C\xrightarrow{\chi_{B,C}}(GB)^{GC}
                \xrightarrow{\tau^\kappa}B^C.
\end{equation}
\end{enumerate}
\end{proposition}

\begin{proof}
For independent variables $A,B,C$, closedness and the lax tensor map give
\begin{align*}
 \Map(A,B^C)
 &\simeq\Map(A\otimes B,C)\\
 &\longrightarrow\Map(G(A\otimes B),GC)\\
 &\longrightarrow\Map(GA\otimes GB,GC)\\
 &\simeq\Map(GA,(GB)^{GC}).
\end{align*}
The first middle arrow is the map on mapping spaces induced by the $\infty$-functor $G$; the second
precomposes with $\mu_{A,B}$.
This composite is natural on
$\cC^{\op}\times\cC^{\op}\times\cC$.

Write $y$ for the Yoneda embedding and $G^*$ for precomposition by $G$ on presheaves, enlarging the universe if necessary.
The above composite is a natural transformation
\[
 yB^C\longrightarrow G^*y (GB)^{GC}
\]
of presheaves in the independent variable $A$.
This is a transformation of presheaf-valued functors of $(B,C)$.
The functor $\infty$-category of presheaves is presentable, and precomposition $G^*$ preserves limits (they are computed pointwise), so $G^*$ has a left adjoint $G_!$ by the adjoint functor theorem \cite[Corollary 5.5.2.9]{HTT}.
For $P$ a presheaf, the adjunction and the Yoneda lemma give
$\Map(G_!y(X),P)\simeq (G^*P)(X)=P(GX)\simeq\Map(y(GX),P)$,
hence $G_!y\simeq yG$
\cite[Proposition 5.1.3.1]{HTT}.
Taking the mate in the functor category with parameters $(B,C)$
therefore gives a natural transformation
\[
 yG B^C\longrightarrow y (GB)^{GC}.
\]
Full faithfulness of the Yoneda embedding yields the natural transformation $\chi$.
Under the tensor-Hom adjunction, $\chi_{B,C}$ corresponds to the composite
\[
 G B^C\otimes GB\xrightarrow{\mu}
 G(B^C\otimes B)\xrightarrow{G(\ev)}GC.
\]

On $\Frob(\cC,G)^{\op}\times\Cart(\cC,G)$, the universal structure
arrows $\tau$ and $\kappa$ in the defining pullbacks are natural
transformations.
Contravariance in the first variable and covariance in the second make
$\tau^\kappa$ a natural transformation
$(GB)^{GC}\to B^C$ on this product.
Thus \eqref{eq:abstract-hom-structure} is a natural transformation $G(B^C)\to B^C$ on that product of lax equalizers.
The pullback universal property of $\Cart(\cC,G)$ gives the claimed functor.
\end{proof}

\begin{thebibliography}{99}
\bibitem[HA]{HA} J.~Lurie, \emph{Higher Algebra}, version of 18 September 2017, \url{https://www.math.ias.edu/~lurie/papers/HA.pdf}.
\bibitem[HTT]{HTT} J.~Lurie, \emph{Higher Topos Theory}, Annals of Mathematics Studies 170, Princeton University Press, 2009.
\end{thebibliography}
\end{document}
